\section{Model}
\label{sec:model}

\subsection{Three POMP models}

Experiment~4 uses one data-generating POMP and two fitting POMPs
(Figure~\ref{fig:model-hierarchy}).  All three models contain the same
stochastic epidemic transitions and negative-binomial observation process.
They differ only in the source of the transmission rate used by the epidemic
process.  The data-generating model uses a prescribed deterministic path
\(\Btrue(t)\); the Gamma-noise fitting model includes \(B(t)\) as a positive
latent stochastic state; and the constant-\(B\) fitting model uses one positive
static parameter \(\Bconst\).

This separation makes the comparison interpretable.  The generator defines
the target against which recovery is measured, whereas the two fitting models
define alternative explanations of the same observations.  Calling the
prescribed path deterministic refers only to its transmission-rate input: the
epidemic trajectory and the reported cases generated from it are random.  In
the same way, calling the comparator constant refers only to its transmission
rate; its epidemic and observation processes remain stochastic.

\begin{figure}[htbp]
\centering
\begin{tikzpicture}[
  node distance=0.8cm and 1.0cm,
  box/.style={
    draw,
    rounded corners,
    align=center,
    minimum height=0.9cm,
    text width=4.0cm,
    inner sep=6pt,
    font=\small
  },
  fitbox/.style={
    draw,
    rounded corners,
    align=center,
    minimum height=0.9cm,
    text width=3.6cm,
    inner sep=6pt,
    font=\small
  },
  arrow/.style={-{Latex[length=2.4mm]},thick}
]
\node[box] (generator)
  {Prescribed \(B_{\mathrm{true}}(t)\)\\data-generating POMP};
\node[box,below=of generator] (data)
  {One reported-case series \(Y_{1:N}\)};
\node[fitbox,below left=0.9cm and 0.4cm of data] (gamma)
  {Gamma-noise \(B(t)\)\\fitting POMP};
\node[fitbox,below right=0.9cm and 0.4cm of data] (constant)
  {Constant-\(B\)\\fitting POMP};
\draw[arrow] (generator) -- (data);
\draw[arrow] (data) -- (gamma);
\draw[arrow] (data) -- (constant);
\end{tikzpicture}
\caption{Three-model design.  Each simulated case series is supplied unchanged
to both fitting POMPs.}
\label{fig:model-hierarchy}
\end{figure}

Immediately before measurement at observation time \(t_n\), the Gamma-noise
state is
\[
X_n^{(G)}
=\{S(t_n),I(t_n),R(t_n),H(t_n),B(t_n)\}.
\]
The data-generating and constant-\(B\) states omit the dynamic transmission
component:
\[
X_n^{(D)}=X_n^{(C)}
=\{S(t_n),I(t_n),R(t_n),H(t_n)\}.
\]
The variable \(H\) is an incidence accumulator.  It records infections since
the previous observation and is reset after the corresponding measurement is
generated.

\subsection{Stochastic SIR process}

Let \(S_k\), \(I_k\), and \(R_k\) be the susceptible, infectious, and recovered
counts at Euler time \(t_k\), with step length \(\Delta t\).  Write
\(B_k^{\mathrm{use}}\) for the transmission value supplied by the relevant
model.  The one-step infection and recovery probabilities are
\begin{align}
p_{SI,k}
&=1-\exp\left(-B_k^{\mathrm{use}}I_k\Delta t/N\right),&
p_{IR,k}
&=1-\exp\left(-\mu_{IR}\Delta t\right).
\end{align}
Conditional on the current state, the event counts are
\begin{align}
\Delta N_{SI,k}
&\sim\operatorname{Binomial}(S_k,p_{SI,k}),&
\Delta N_{IR,k}
&\sim\operatorname{Binomial}(I_k,p_{IR,k}).
\end{align}
The factor \(B_k^{\mathrm{use}}I_k/N\) is the per-susceptible infection
hazard.  The exponential conversion gives the probability of infection over
the finite step \(\Delta t\), while the recovery probability is derived from
the constant recovery hazard \(\mu_{IR}\).  This construction keeps both
probabilities in \([0,1]\) and lets the same transition equations accept any
of the three transmission inputs.
The state then changes according to
\begin{align}
S_{k+1} &= S_k-\Delta N_{SI,k},\\
I_{k+1} &= I_k+\Delta N_{SI,k}-\Delta N_{IR,k},\\
R_{k+1} &= R_k+\Delta N_{IR,k},\\
H_{k+1} &= H_k+\Delta N_{SI,k}.
\end{align}
Thus \(S_k+I_k+R_k=N\), while \(H_k\) records events rather than people and is
excluded from the population total.  The binomial draws make the epidemic
process stochastic even when \(B_k^{\mathrm{use}}\) is deterministic or
constant.

The accumulator bridges the process and observation time scales.  The SIR
state is updated every Euler step, but cases are observed only once per day.
Summing \(\Delta N_{SI,k}\) into \(H\) preserves all infections between two
measurements; resetting \(H\) after measurement prevents the same infection
from contributing to more than one reported count.

\subsection{Observation model}

Let \(H_n\) be the infections accumulated during
\((t_{n-1},t_n]\), immediately before reset, and let \(Y_n\) be the reported
case count.  Reported cases depend on incidence \(H_n\), not on the number
infectious at one instant.  The model is
\begin{equation}
Y_n\mid H_n
\sim
\operatorname{NegBin}
\left(\text{mean}=\rho H_n,\ \text{size}=k\right),
\label{eq:measurement}
\end{equation}
with
\[
\operatorname{E}(Y_n\mid H_n)=\rho H_n,
\qquad
\operatorname{Var}(Y_n\mid H_n)
=\rho H_n+\frac{(\rho H_n)^2}{k}.
\]
Here \(\rho\) controls the expected reported fraction and \(k\) controls
overdispersion around that conditional mean.  The observation model does not
feed reports back into the epidemic: \(Y_n\) is generated from \(H_n\), after
which the latent SIR process continues from its current state.
This negative-binomial variation is observation noise conditional on latent
incidence.  It is distinct from the binomial demographic variation in the SIR
transitions.

\subsection{Transmission-rate models}

\paragraph{Gamma-noise model.}
The Gamma-noise fitting POMP initializes \(B(0)=B_0>0\).  At each Euler step it
draws
\begin{equation}
B_{k+1}\mid B_k
\sim
\operatorname{Gamma}
\left(
  \text{shape}=\frac{1}{\sigB^2\Delta t},
  \ \text{scale}=B_k\sigB^2\Delta t
\right),
\label{eq:gamma-update}
\end{equation}
and uses the updated value \(B_{k+1}\) immediately in the infection
probability.  The two static transmission parameters are \(B_0\) and
\(\sigB\); the sequence \(B(t)\) remains latent.

For the shape--scale parameterization in
Equation~\ref{eq:gamma-update},
\begin{align}
\operatorname{E}(B_{k+1}\mid B_k)
&=B_k,\\
\operatorname{Var}(B_{k+1}\mid B_k)
&=B_k^2\sigB^2\Delta t.
\end{align}
The transition therefore has zero conditional drift and one-step conditional
coefficient of variation \(\sigB\sqrt{\Delta t}\).  Larger \(\sigB\) permits
greater relative movement between adjacent process times.  Because the Gamma
distribution has positive support, \(B(t)\) remains positive.  When time is
measured in weeks, \(B(t)\) has units \(\mathrm{week}^{-1}\) and \(\sigB\)
has units \(\mathrm{week}^{-1/2}\).

Zero conditional drift does not force a realized path to remain flat.  A draw
can move above or below the current value, and successive innovations can
accumulate into a sustained change.  Conversely, the transition contains no
preferred change point, downward trend, or mean-reversion level.  Evidence for
a decline must therefore enter through the observed cases during filtering.
The role of \(\sigB\) is specifically to regulate local relative process
variation; it is not a reporting-dispersion parameter and it is not an IF2
search perturbation.

\paragraph{Prescribed data-generating rate.}
The data-generating POMP sets
\[
B_k^{\mathrm{use}}=\Btrue(t_k).
\]
This path is fixed by the simulation design rather than drawn from
Equation~\ref{eq:gamma-update}.  The generated epidemic and reports are still
random because their binomial transitions and negative-binomial measurements
remain stochastic.  Section~\ref{sec:design} gives the prescribed path and
parameter values.

\paragraph{Constant-\(B\) model.}
The constant fitting POMP estimates one positive parameter \(\Bconst\) and
sets
\[
B_k^{\mathrm{use}}=\Bconst
\]
throughout the epidemic.  It has no latent transmission-rate state, but it
retains the same stochastic SIR and observation processes as the other two
POMPs.

The scientific dependence common to all three models is
\[
B(t)
\longrightarrow
\text{infection hazard}
\longrightarrow
\text{new infections}
\longrightarrow
H_n
\longrightarrow
Y_n.
\]
The Gamma-noise model adds stochastic variation at the first link.  Particle
filtering and IF2, introduced next, are computational tools for inference and
do not add scientific transitions to this chain.

Five sources of variation should therefore be kept separate.  Gamma process
noise moves the latent transmission rate only in the Gamma-noise fitting POMP;
binomial demographic noise moves the epidemic counts in all three POMPs; and
negative-binomial observation noise moves reported cases around latent
incidence.  Particle-filter Monte Carlo error is numerical approximation error
from using finitely many particles, while IF2 perturbations are temporary
optimization devices.  The last two affect inference but are not components
of the scientific epidemic model.
